\section{Introduction}\label{sec:intro}
It is now routine to improve an LLM system by searching over prompts or multi-call programs against a benchmark---by hand, with random or grid search, or with compilers such as DSPy \citep{khattab2023dspy,opsahl2024mipro}---and then to report the best configuration's score on that same benchmark.
The benchmark is then used twice: to \emph{choose} the configuration and to \emph{score} it.
Because prompt choices move accuracy by several points while differing only slightly in true quality \citep{sclar2024quantifying}, the maximum of $K$ noisy estimates on the same questions is optimistic, much as in post-selection inference \citep{berk2013valid,andrews2024inference} and adaptive data analysis \citep{dwork2015preserving,blum2015ladder}.
What is unclear is (i) how large this gap is for realistic LLM prompt search and what drives it, (ii) whether it matters for the decisions that reported numbers feed---claims of improvement, comparisons between systems with different tuning effort, deployment beyond the tuning distribution---and (iii) whether there is a reporting practice that is reliable, and how one would tell.

\textbf{We study these questions with a theory-plus-experiment design} (full-benchmark correctness matrices cached from real LLMs on MMLU-Pro, resampled to simulate $10^3$ tuning episodes per setting, with the true accuracy of the delivered configuration known exactly, plus a synthetic generator and a check on fresh questions).
Our findings, all within the (limited) scope of Sec.~\ref{sec:limits}:
\begin{itemize}[leftmargin=1.2em,itemsep=0pt,topsep=2pt]
\item \textbf{The gap is large and predictable.} With 16--32 candidate prompts and a 100-question benchmark, the reported accuracy of the delivered prompt exceeds its true accuracy by 4--5 points (Qwen2.5-0.5B/1.5B, MMLU-Pro), 2 points at $n{=}400$; it follows $\sigma\sqrt{(1-\rho)/n}\cdot\mathbb E[\max_K]$ where $\rho$ is the correlation between candidates' per-question outcomes, and it vanishes only when one candidate is clearly best. At $n\le400$ at most half of the reported gain over an average candidate is real, and the reported gain always exceeds that of the best candidate (Sec.~\ref{sec:magnitude}).
\item \textbf{It changes conclusions.} Against the best prompt as baseline (no candidate is truly better), naive significance tests declare an improvement in 12--15\% of runs instead of 5\%; a system tuned over 16 candidates looks $+4.0$ points better than an equally good one tuned over 2 (true: $+0.7$). In-distribution protocols do not address a separate subject shift (Sec.~\ref{sec:decisions}).
\item \textbf{A simple, reliable practice exists, with a measurable price.} Selecting on a small share (e.g.\ 20\%) of the questions and reporting the held-out remainder is unbiased with nominal coverage for \emph{any} search procedure, and at $n\le800$ has lower RMSE than the naive number; max-$T$ simultaneous bounds are valid for any selection rule over the candidate set but conservative; nested CV is an unbiased point estimate but its Wilson interval under-covers; bootstrap-optimism correction and empirical-Bayes shrinkage are \emph{not} reliable (residual bias; failure when one candidate is truly better) (Sec.~\ref{sec:protocols_results}). A protocol should be validated, as here, by simulating the whole search--report pipeline against known truth and, as a check of the simulation, on fresh questions (Sec.~\ref{sec:fresh}).
\end{itemize}
We do not propose new estimators; the contribution is the quantification for this setting, the decision-level consequences, and the comparison of remedies, including negative results.
